\chapter{Components worth testing}
\label{ch:software-dossiers}

Imagine a two-page research note. Its opening reports an estimated premium of
42 basis points. A table displays the same estimate, an equation defines the
premium, and a citation supports the identification strategy. The author wants
help with one awkward paragraph.

A conventional writing assistant sees the paragraph. Humanvoice must see the
note. If a proposed edit changes 42 to 24, removes the condition attached to the
estimate, or moves the citation to a stronger claim, the new prose may still be
fluent. The revision has nevertheless changed the analysis. The product needs
to show the author that change before anyone accepts the wording.

This small note gives the software review its order. First preserve the source
and build what the reader sees. Then locate the prose question. Compare the
meaning-bearing parts of the two versions. Only after those steps should an
author decide whether to revise and a reader decide whether the note works.
Available packages can perform many of these jobs, but no package in the survey
performs the complete sequence.

The detailed component and market records appear in
Appendix~\ref{app:component-records}. This chapter asks the narrower investment
question: which parts should the first build reuse, which need an adapter, and
which connection must Humanvoice own?

\section{One note, six jobs}
\label{sec:component-jobs}

The note first needs an identity. Humanvoice takes a snapshot of the LaTeX
source, its bibliography, and the command that produces the PDF. A later
finding can then point to the exact version on which it was made. Without that
identity, an author may open a warning from yesterday against a paragraph
changed this morning.

Next comes the build. The software compiles the note and records unresolved
references, missing files, and the page on which each labelled object appears.
A source file that looks plausible but does not produce the reader's PDF is not
a useful basis for revision.

The third job is diagnosis. A prose rule may notice that the paragraph begins
with a vague statement of importance. A structural check may notice that the
same opening appears in four sections. A citation check may notice that the key
no longer resolves to the source used in the previous version. Each finding
needs a location and a reason. None changes the source by itself.

The fourth job begins only after the author edits. Humanvoice compares the old
and new versions of the equation, the number, the citation identity, and the
qualification around the result. In the example, the author should see that the
premium remains 42 basis points in the paragraph, table, and equation. If one
instance differs, the comparison remains open.

Privacy supplies a fifth job. A restricted manuscript stays local unless its
owner approves a particular external service and the fields that may be sent.
Installing a local model can reduce one class of disclosure risk. It does not
make the model's suggestion correct, and it does not erase the need to record
which version and prompt produced the suggestion.

The last job belongs to a person who did not write the note. That reader sees
the rendered pages, the declared decision, and any required disclosure. Can the
reader state what the 42 basis points measure, why the estimate is identified,
and what action follows? The answer is the product outcome. It cannot be
recovered from a grammar score or a successful build.

No current candidate does all six jobs. That is not a reason to write six new
systems. It is a reason to put mature components behind a small record that
keeps their evidence separate. A parser supplies locations. A linter supplies
warnings. A model supplies suggestions. A reader supplies the consequential
judgment.

The path from an interesting repository to an adopted component should be
deliberately slow.

\hvFigurePackageFunnel

A package that installs has shown only that it can run. A package that handles
one designed example has shown a little more. It earns a role in Humanvoice
only after it succeeds and fails in the expected places on examples it did not
help to design. Even then, adoption is conditional on the burden it creates for
the author.

\section{Start with the source, not the suggestion}
\label{sec:component-source}

LaTeX is a program as well as a document format. A command can define a number
once and print it in several places. Another command can alter spacing without
altering meaning. A third can expand differently inside a table and a sentence.
The parser therefore has to preserve both the raw source and what the command
produces.

\begin{cardexample}{One command, two visible meanings}
Consider a fixture built around a custom command called
\verb|\riskterm|. In the table it expands to ``42 bp.'' In the prose it expands
to ``estimated risk premium.'' An editor that stores only the expanded text
sees two unrelated phrases. An editor that stores only the command name misses
the visible difference. Humanvoice needs the call, the expansion, the source
location, and the rendered location. If any of those cannot be recovered, the
fixture should be marked unparsed rather than protected.
\end{cardexample}

pylatexenc is the smallest serious Python candidate for this first pass
\citep{pylatexenc2026}. It can tokenize commands and groups while retaining
source positions. That is enough to test whether the product can locate known
macros, mathematics, citations, and environments without inventing a full
model of TeX. Its limitation is equally useful: an unfamiliar macro has no
semantic meaning merely because its braces were parsed. The adapter must keep
the raw span and abstain.

\Needspace{13cm}
A richer document model may become necessary. plasTeX understands sections,
counters, labels, and more of the structure a rendered document exposes
\citep{plastex2026}. LaTeXML translates many mathematical constructs into a
semantic XML representation \citep{latexml2026}. Both can answer questions a
tokenizer cannot. Both also create a larger compatibility problem. A custom
package may behave differently from the author's TeX installation. The
benchmark must compare their representation with the canonical PDF rather than
assuming that a more elaborate tree is more faithful.

The editor route creates a different tradeoff. unified-latex offers typed
JavaScript structures and transformations that fit a web interface
\citep{unifiedlatex2026}. tree-sitter-latex can update a syntax tree quickly as
the author types \citep{treesitterlatex2026}. TexSoup offers a tolerant Python
interface that remains usable on imperfect input \citep{texsoup2026}. Speed and
tolerance are helpful in an editor. They become dangerous when a best-effort
parse is silently promoted to a claim that an equation was protected. Unknown
constructs must remain visible in every route.

Pandoc solves another job well. It can turn supported sources into a structured
interchange form and produce a prose view for linguistic diagnostics
\citep{pandoc2026filters}. That secondary view is useful when raw commands
confuse a writing rule. It is not the authority for equation preservation. The
LaTeX source and its own build remain the reference for the first release.

latexdiff belongs at the end of this path, not the beginning. It gives an author
a legible picture of source changes \citep{latexdiff2026}. That picture is
valuable in a review packet. It does not decide whether two equations are
equivalent or whether a citation still supports a sentence. Humanvoice should
use it to display an already recorded comparison, never to define the
comparison.

The first parser decision is therefore a bake-off, not a commitment to a
framework. Give the candidates the same custom commands, nested environments,
equations, tables, comments, and malformed input. Record what each locates,
what it loses, and where it abstains. Begin with pylatexenc because it is small
and fits the proposed Python core. Keep the richer candidates available when a
fixture demonstrates that the smaller representation is insufficient.

\section{Let mature tools do narrow jobs}
\label{sec:component-narrow-tools}

The product does not need to conceal the ordinary LaTeX toolchain. latexmk can
repeat compilation until references and the bibliography settle
\citep{latexmk2026}. ChkTeX can report source-level syntax and usage warnings
\citep{chktex2026}. Their output belongs in the run record with the exact build
command. Neither tool needs to become a Humanvoice feature, and neither tool
establishes that the resulting page is understandable.

Grammar and prose checks are another mature layer. Vale is well suited to
project rules because it returns a line, rule, and message
\citep{vale2026docs}. LTeX+ and TeXtidote understand enough LaTeX to avoid
treating every equation as a sentence and can return language findings to
source locations \citep{ltexplus2026}. textlint is attractive when a rule needs
a programmable JavaScript host \citep{textlint2026docs}. The first build should
adapt their machine-readable findings rather than copy their rule engines.

\hvFigureNarrowEvidence

These tools should run after the important question is visible. If the research
note does not explain what the 42 basis points measure, a queue of comma and
agreement corrections is a distraction. Once the mechanism is present, those
corrections can reduce friction. The reading brief determines the order; the
number of warnings does not.

\Needspace{13cm}
Citations require two separate checks. The first resolves identity: does the
key still point to the expected authors, title, venue, and identifier? GROBID
can help extract structured fields when a PDF is the only available source
\citep{grobid2026}. The second check asks whether the paper supports the
sentence. Matching a DOI cannot answer that question. Humanvoice should never
label an identity match as citation validation.

Local inference is also a component rather than a conclusion. llama.cpp offers
a portable runtime for compatible local models \citep{llamacpp2026}. It can
keep an approved prompt and manuscript on the organization's hardware. The
selected model still needs a licence, resource profile, prompt-injection test,
and comparison with a remote baseline. A local hallucination is private, but it
is still a hallucination.

The reader task needs less machinery than a general model-evaluation platform.
jsPsych can present a local browser exercise, randomize version order, record
answers and elapsed time, and export a structured session
\citep{jspsych2026}. That makes it a plausible instrument for the first human
comparison. It does not recruit an expert reader or decide whether an answer is
correct.

Inspect AI can organize repeatable tasks, solvers, and model graders
\citep{inspectai2026}. It may become useful when the team compares prompts or
judge models. Making it a dependency of the first human exercise would reverse
the order of evidence. The named reader is the endpoint; model-evaluation
infrastructure is optional support.

\subsection{The pacing adapter is a small local experiment}
\label{sec:component-pacing-adapter}

The HPO case adds one capability that the general packages do not supply. The
first build needs a way to compare a target chapter with the exemplars in its
reading brief and to show where new concepts arrive in bursts. The repository
now contains a standard-library prototype, \texttt{tools/concept\_density.py},
for that narrow job. It recognizes a few transparent signals---acronyms,
mid-sentence proper terms, emphasized definition phrases, recurring technical
bigrams, and citation keys---then reports introduction rate, consolidation,
bursts, and live load.

The prototype is deliberately not a dependency of the protected comparison.
It can be noisy, and it has no sentence-level dependency parser. Its output is
a ranked worklist with file and line locations. The author can dismiss a
capitalized ordinary word in one action. A first-use pass reads the same brief
for known vocabulary and exemption ranges, then asks for a gloss or a
taught-later pointer when the evidence is missing.

The HPO fixture supplies the first calibration and held-out material: 25
adjudicated ordering violations and 80 adjudicated false alarms. Those numbers
remain in the fixture record rather than becoming a success rate. Before the
adapter is promoted, an independent operator must measure precision,
abstention, and minutes spent triaging it on documents outside its tuning
partition. A positive result would justify a narrow adapter; a null result
would leave the rest of the source-to-reader path intact.

The common rule is simple. Reuse a mature tool when its output names a fact the
product can check: a build result, a source location, an identifier, a language
warning, a runtime trace, or a recorded response. Do not ask that tool to make
the larger claim that the document is correct or persuasive.

\section{The alternatives an executive will try first}
\label{sec:component-alternatives}

A sponsor should ask why the team cannot buy this capability. The answer cannot
be that Humanvoice has a longer feature list. Existing assistants are better
than a new prototype at many common writing tasks, and the experiment should
retain them as credible alternatives.

For an ordinary memorandum, Grammarly may be the right incumbent. Its
enterprise product works across common writing environments and combines
language suggestions with a broader agent layer
\citep{grammarly2026enterprise}. Its privacy documentation also explains when
generative features may send text and context to service providers
\citep{grammarly2026privacy}. An organization can weigh those controls against
the document's sensitivity. Humanvoice should not spend its first budget
rebuilding general grammar and tone advice that a mature service already
provides.

The research-note example exposes the remaining question. The unit seen by a
general assistant is usually a selected passage or communication context. The
unit Humanvoice must protect is the versioned note: paragraph, equation, table,
citation, and reader decision together. An excellent rewrite of the paragraph
can still change the denominator or remove the identification condition.
Generic fluency is therefore part of the incumbent workflow, not the proposed
product wedge.

\hvFigureAlternativeUnits

\Needspace{13cm}
Academic assistants come closer. Writefull integrates research-trained language
support into Overleaf and offers rewriting and style functions in the LaTeX
environment \citep{writefull2026overleaf}. It should be tested as a language
adapter and as an incumbent arm. The question is not whether it can improve a
sentence. The fixture asks whether the suggestion remains tied to a source
span, whether protected objects stay visible, and whether the resulting note
can enter the same reader task.

Paperpal combines language editing with reference discovery, citation styles,
and submission checks \citep{paperpal2026assistant}. Trinka similarly targets
academic and technical writing and documents institutional, privacy,
integration, and API options \citep{trinka2026privacy}. Their breadth is real
value for a researcher. It also makes causal interpretation harder: when one
interface performs many interventions, a favorable reaction does not show
which intervention helped or whether a protected object moved. The benchmark
should test declared functions one at a time.

Overleaf already solves the social side of collaborative LaTeX editing. It
provides comments, version history, and Track Changes through which authors can
accept or reject revisions \citep{overleaf2026trackchanges}. Humanvoice should
integrate with that environment when a team already uses it. Track Changes is a
view of source edits, however, not a proof that an equation is equivalent or a
citation supports its new sentence. The source snapshot and protected
comparison remain distinct.

Privacy does not create an automatic winner. A hosted assistant may be suitable
for a public draft and unacceptable for a restricted manuscript. A local model
may satisfy the data boundary and perform poorly on the task. An editor may
retain complete version history while the author prefers not to expose that
history to an external reader. Each workflow should record what leaves the
machine, what returns, and what the reader receives.

The fair comparison is therefore practical. Give the incumbent and Humanvoice
the same note, reading brief, and deliberate defects. Ask each route to surface
a changed number, a lost qualification, and a moved citation. Record network
calls, author time, false passes, and the reader's ability to reconstruct the
result. If an existing product performs the complete job at lower burden, the
team should adopt it and narrow or stop Humanvoice.

\section{Buy the parts, build the connection}
\label{sec:component-build-boundary}

The software survey leads to a restrained boundary. Buy or reuse grammar,
collaborative editing, and document compilation when the licence and privacy
terms fit. Adapt parsers, linters, and comparison views when they can return a
stable location and an inspectable reason. Build the connection that current
alternatives leave implicit.

That connection begins with a source snapshot. It records the note, build, and
reading brief as one run. A source map then links the paragraph, equation,
number, table, citation, and rendered page. Diagnostics attach findings to
those locations. The author accepts, rejects, or defers each finding. A new run
records the revision and compares its protected parts with the old one. The
reader's response closes the sequence.

Return to the awkward paragraph. Vale may flag its vague opening. The author
rewrites it in Overleaf. Humanvoice does not need to own either operation. It
does need to know which source version produced the finding, which version
contains the repair, and whether 42 basis points and the identification
condition survived. It also needs to place the revised PDF in front of the
reader without exposing private comments or model prompts.

Five small records are sufficient for the first experiment: the run, protected
object, finding, author decision, and reader response. Immutable identifiers
join them. This is not a general document platform. It is the minimum structure
needed to answer who saw what, what changed, and which observation supports the
next investment decision.

The boundary affects cost. Reimplementing a grammar engine would consume the
MVP on a solved problem. Depending entirely on a cloud editor would save
engineering time while moving the central source comparison into a service we
do not control. The pilot should therefore price three configurations: a fully
local route, an incumbent-editor route with Humanvoice adapters, and a hosted
academic assistant used only on approved excerpts.

The choice among those configurations belongs to the experiment. The local
route may cost more operator time. The editor route may provide the best author
experience. The hosted route may produce the best language suggestions. Reader
performance, meaning errors, privacy compliance, and total burden decide; the
architecture diagram does not.

\clearpage
\section{A fixture decides, not a feature list}
\label{sec:component-fixture}

The first held-out set should be small enough to understand completely. One
fixture is clean. One contains a private project label. One changes an exponent
in an equation. One changes the denominator in a table. One cites a real paper
that does not support the sentence. One uses the custom \verb|\riskterm|
command. One contains malformed LaTeX. One repeats a sentence shape on purpose.

Each fixture states what success and failure look like before a package runs.
For the custom command, a parser succeeds if it preserves the call, expansion,
and locations. It abstains acceptably if it identifies the command as unknown
and blocks a protected pass. It fails dangerously if it returns a complete-
looking record that has lost either use of the command.

For candidate $k$ on fixture $i$, record whether the capability was observed,
whether a false pass occurred, how long the run took, and how much operator
work was required. Keep those observations as a matrix. A single score would
allow fast execution to compensate for silent loss of an equation, which is
not an admissible trade.

The benchmark also records ordinary engineering facts. A restrictive licence
may confine a useful package to a reference implementation. A stale release may
be acceptable for a frozen migration and unsuitable for a service that must
track new TeX packages. A component that contacts the network may be excluded
from a private note even when its findings are good. These facts sit beside the
functional result rather than being folded into a popularity ranking.

Adoption requires three observations. The component must find the positive
case. It must fail or abstain on the negative case. A second operator must be
able to replay the finding from the recorded version and configuration. A
package that cannot meet all three may remain useful to study, but it does not
enter the protected path.

The same standard applies to the product as a whole. If the components cannot
produce a stable source map, the project remains a prose linter. If the source
map works but authors face an unmanageable queue, the workflow must be reduced.
If the engineering succeeds but the named reader does no better, the proposed
connection has not earned its cost.

This component strategy leaves the first build with one narrow responsibility:
join reliable observations without inflating what any of them proves. The next
part follows a document through that connection from draft to reader decision.
